\chapter{Conclusions and Future Work}
\label{sec:conclusion}
\lhead{\emph{Chapter 5. Conclusions and Future Work}}  % page header

This report set out to test whether research in machine learning herds. It separated two ways in which a field can
conform, imitation of what others are doing and continuity with one's own record, and tested their predictions on ten
years of ICLR submissions at three levels: the field as a whole, entry to a topic and the acceptance decision. This
chapter first summarises the key findings (Section~\ref{sec:findings}). It then discusses what they mean for each
research question and how they relate to earlier work (Section~\ref{sec:discussion}), sets out the limitations
(Section~\ref{sec:limitations}) and suggests future work (Section~\ref{sec:future}).

\section{Summary of key findings}
\label{sec:findings}

At the level of the field (RQ1), keyword waves are real. Of the 68 most frequent author keywords, 17 rose and then fell
between 2019 and 2026 (\lab{exploratory}). The number of topics that appear to be emerging, however, depends on how the
topics are defined. The same papers give 8 emerging topics when the topic map is fitted on 2017--2020 and 29 when it is
fitted on all years (\lab{exploratory}). The one change that clearly lasted is the shift to foundation models: their
share of submissions rose from 0.3\% in 2017 to 54.4\% in 2026, and topics in which more papers built on FMs in the
previous year have fewer papers proposing a new architecture, by 8.0\,pp in discovery and 11.0\,pp on held-out papers
for a change from no adoption to full adoption (\lab{replicated}). The registered tests of whether a topic's mix of
contribution types changes as the topic grows were not supported.

At entry to a topic (RQ2), the registered herding prediction failed in the opposite direction. Papers by newcomers to a
topic are more often method papers than papers by incumbents, by 3.7\,pp in discovery and 3.6\,pp on held-out papers,
and less often measurement papers, by 1.7\,pp and 2.1\,pp (\lab{replicated}). The method surplus is linked to entering
the topic rather than to being new to ICLR, since researchers who already had ICLR papers in other topics show it as
well, and it appears in every period, including the years before 2023 (\lab{exploratory}). Much of the measurement gap
goes with team size, and in 2024--2026 measurement papers are concentrated among incumbents and larger teams. Newcomer
papers also lie farther from earlier ICLR work, and first-time ICLR teams take up foundation models and young, rising
keywords later than teams with an ICLR record (\lab{exploratory}). Whether entrants flock to hot topics could not be
decided (\lab{inconclusive}).

At the acceptance decision (RQ3), the registered test found no detectable penalty or premium for papers in crowded
topics, but the test could not rule out effects smaller than about 5\,pp per log unit (\lab{uninformative}). Branded
titles spread from 10.9\% to 44.6\% of titles and carry a small premium in discovery, which the held-out test was too
weak to confirm (\lab{exploratory}). The associations that held up on held-out papers all involve the team's own
earlier work. Papers far from it are accepted 5.4\,pp less often in discovery and 6.1\,pp less often on held-out
papers; for the same author, a paper closer to their earlier work is accepted more often; and a team with one topic
insider is accepted more often than a team with none (all \lab{held-out supported}, although the insider result does
not survive a family-wise correction). A topic that was favoured at review in 2018--2022 kept only part of its
advantage in 2024--2026 (\lab{held-out supported}). FM papers are accepted more often, by 4.6\,pp, and this premium did
not shrink as FMs spread within a topic, while the timing of keywords shows at most weak associations with acceptance
(both \lab{exploratory}).

\section{Discussion}
\label{sec:discussion}

\subsection{Does machine learning herd?}

Table~\ref{tab:routes} places each result under the route it fits. The field converged on foundation models, but
imitation, continuity, lock-in and efficient choice all predict that, so the field level cannot settle the question.
Where the two routes make different predictions, at entry and at the gate, the evidence leans towards continuity. The
registered herding predictions for entry failed or were inconclusive, and the registered crowding test at the gate was
uninformative. Every result that replicated at the gate is referenced to the team's own record, with one exception:
the persistence of topic edges, which fits imitation of past acceptance and regression to the mean in a noisy review
process equally well. The conformity that can be detected at ICLR is a pull towards a record rather than towards what
others are doing.

\begin{table}[tb]
\centering
\footnotesize
\setlength{\tabcolsep}{3pt}
\renewcommand{\arraystretch}{1.15}
\begin{tabularx}{\textwidth}{@{}L{2.3cm} Y Y@{}}
\toprule
 & Entry (what is written) & Gate (acceptance) \\
\midrule
\multicolumn{3}{@{}p{\textwidth}@{}}{\textbf{Field}, both routes: convergence on foundation models (\lab{replicated}),
also predicted by lock-in and efficient choice.} \\
\midrule
\textbf{Imitation} (others, now) & Follow-on prediction reversed (\lab{replicated}); newcomers farther from earlier
work and fewer twins for returning topic newcomers (\lab{exploratory}); flocking \lab{inconclusive}. & No detectable
crowding premium (\lab{uninformative}); branded titles: small premium, held-out test too weak (\lab{exploratory}).
Also fit a rival: topic edges partly kept (\lab{held-out supported}); no change in the FM premium; weak keyword timing
(\lab{exploratory}). \\
\midrule
\textbf{Continuity} (a record) & Incumbents write more measurement papers (\lab{replicated}; much of it with team size);
foundation models spread through teams with an ICLR record, as diffusion also predicts (\lab{exploratory}). & Lower
acceptance far from the team's record; higher for the same author's closer work and with a topic insider
(\lab{held-out supported}; status and expertise fit too). \\
\bottomrule
\end{tabularx}
\caption[Results by route and level]{Where conformity shows, by route and level. The table summarises
Chapter~\ref{sec:results} and adds no new estimates.}
\label{tab:routes}
\end{table}

This answer needs two qualifications. First, the absence of a detectable crowding premium does not rule out a
bandwagon. \citet{abrahamson1993institutional} showed that bandwagons can spread even when adopting brings no gain, and
the spread of branded titles may be such a case. Second, the record-side results at the gate fit continuity, but they
do not single it out. The next subsections take each research question in turn.

\subsection{The field: trends and hindsight}

The hindsight result speaks to a common practice in the history of research fields. Studies of PNAS, computational
linguistics and AI research have fitted one topic model to all years and then read trends from it
\citep{griffiths2004finding,hall2008studying,anderson2012towards,klinger2020narrowing}. The present results show that,
with the papers and the pipeline held fixed, the choice of fitting years alone can change the count of emerging topics
by a factor of more than three. This turns a concern raised in principle by \citet{leydesdorff2002dynamic} and
\citet{rotolo2015what} into a measured difference, and it adds the time window to the list of analysis choices that a
multiverse analysis would vary \citep{steegen2016increasing}. It also fits with recent work that tests signals of
emergence using only earlier data \citep{khanbayov2026topical}. The comparison is a bracket, not a correction: the
early map cannot see topics that appeared after 2020, so neither count is the true number.

The shift to foundation models is the clearest case of convergence in the corpus, and it matches the accounts of
paradigm shifts in NLP \citep{gururaja2023build,saphra2024first} and the concern about homogenisation around shared
models \citep{bommasani2021opportunities}. It is also what one would expect if general methods that scale with
computation simply work better \citep{sutton2019bitter}. Because both routes and both of these readings predict it,
the shift by itself says little about herding. This is the reason the report turned to entry and to the gate.

\subsection{Entry: who brings what}

The entry results run against the picture of the newcomer that underlies the herding account. In cascade models, the
person with the least information follows those who chose before them \citep{bikhchandani1992theory,banerjee1992simple}.
At ICLR, researchers new to a topic are the ones who more often propose methods, and their papers lie farther from
earlier work. This fits the literature on outsiders, which finds that people from the margin of a field more often
bring solutions and new ideas \citep{jeppesen2010marginality,zeng2021fresh,azoulay2019does}. The present results add
one detail that earlier work did not measure: the type of contribution. Newcomers bring methods, and incumbents and
larger teams write more of the benchmarks and evaluations. This is consistent with the finding that ML relies on a
small number of benchmark datasets \citep{koch2021reduced}, and with the observation in NLP that tasks and methods lead
while datasets follow \citep{pramanick2023diachronic}. A method paper is not necessarily an original one, however, so
these results say who proposes methods, not whose work is more novel.

First-time teams took up foundation models later than teams with earlier ICLR work. This fits absorptive capacity,
the idea that new knowledge is taken up by those who already hold related knowledge \citep{cohen1990absorptive}. It
fits network diffusion equally well \citep{coleman1957diffusion,rogers2003diffusion}, so it cannot separate the routes.
It does, however, contradict the idea that the newest entrants rush to whatever is fashionable. Half of the first
authors of LLM papers on arXiv in 2023 were new to LLM research \citep{movva2024topics}, but at ICLR, teams new to the
venue did not lead the adoption of foundation models.

\subsection{The gate: crowd or record}

The strongest result at the gate is the far jump: papers far from the team's earlier work are accepted less often.
This extends the pivot penalty, in which larger moves away from one's own past work go with lower impact and lower
publication success \citep{hill2025pivot}, to the accept or reject decision of an ML venue under double-blind review,
with every rejected paper in view. It also fits the finding that innovative research strategies are riskier
\citep{foster2015tradition}. The shape of the association matters: acceptance rises with closeness up to the fourth
fifth of the distribution and then levels off, so leaving one's record is associated with lower acceptance, but staying
very close brings no further gain. This reading suggests a threshold rather than a reward for repetition.

Several explanations fit these results, and the data cannot separate them. Continuity is one. Expertise is another:
reviewers may judge work from a team that knows the area as stronger, and with reason, as grant reviewers are both
biased towards and better informed about work in their own area \citep{li2017expertise}. Status is a third.
Cumulative advantage favours those already recognised \citep{merton1968matthew}, famous authors fare better when
reviewers can see them \citep{tomkins2017reviewer}, and preprints can reveal authors under double-blind review
\citep{bharadhwaj2020deanonymization}; at ICLR, a prestige gap has been tested with a pre-registered out-of-sample test
\citep{ibrahim2026bias}. The insider result is a good example. One insider on the team goes with higher acceptance,
which matches the chaperone effect \citep{sekara2018chaperone} and the finding that senior expertise offsets the
disadvantage of newcomers \citep{biliotti2024breaking}. But each additional insider adds little more than each
additional experienced author from another topic, so a proven record in general may matter more than knowledge of the
specific topic.

On the side of imitation, the gate shows little. The registered crowding test was too weak to detect effects of a few
points, which is consistent with the high level of noise in ML review \citep{cortes2021inconsistency,
beygelzimer2023has}. Topic edges persist only in part, which noise also predicts, although differences between research
areas in how reviewers score papers are a rival account \citep{xu2026reviewer}. Together with a randomised trial that
found no herding among reviewers in discussions \citep{stelmakh2023large}, the results give little support to the idea
that ICLR review rewards following the crowd. Whether evaluators penalise novelty remains contested
\citep{boudreau2016looking,teplitskiy2022is}, and the present results do not settle it: the penalty found here is for
distance from the team's own work, not for distance from the field, and the test of distance from earlier ICLR work as
a whole was uninformative.

\subsection{Why the distinction matters}

The two routes carry different risks for a field. Imitation can discard private information and leave a consensus
that rests on little evidence and breaks easily \citep{bikhchandani1992theory,banerjee1992simple}. Continuity trades
long-run exploration for short-run competence \citep{march1991exploration}, and choices that are safe for each
researcher can make a whole field explore less \citep{rzhetsky2015choosing}. A field that looks herded may in fact
conform by continuity, and a diagnosis of fads would then point to the wrong problem. Acceptance is also only a
short-run outcome. During the previous shift in ML, researchers who stayed close to their own record had lower later
impact \citep{bao2025persistence}, so the work that is accepted and the work that matters later need not be the same. For those who write the history of ML and NLP, the results suggest a simple check: before reading trends from a
topic model, refit it on early years and see which trends remain.

\section{Limitations}
\label{sec:limitations}

All results are associations at one venue, and none measures what would change if a researcher chose a different
topic, team, title or type of contribution. The held-out papers share conference years, reviewers and authors with the
discovery papers. The held-out tests therefore guard against overfitting by the analyst, but not against patterns that
are specific to a year or a pool of reviewers, and a supported held-out test confirms the direction of an association
rather than its size. The standard errors may also be too small: under one simulated null, the main entry tests gave a
significant result 12.6\% and 10.0\% of the time at a nominal 5\% level.

The framework of the report was built partly after the analysis. The distinction between imitation and continuity was
drawn up after the registered and declared tests had run, and the predictions of continuity were written afterwards.
The record-side results at the gate were found post hoc. Two of the three rest on secondary held-out tests, and the
insider result does not survive a family-wise correction. Keyword popularity and title style had been explored against
acceptance before registration.

The measures have their own limits. Contribution and FM labels come from an LLM. They were consistent across two runs,
but they were checked by hand only by the author, on 100 papers, about 20 of which were borderline, and there is no
measure of agreement with a second human coder. A method paper is a type of contribution, not a measure of novelty, and
distance in an embedding space is not originality. Entry status counts ICLR submissions only, and the change in author
identifiers in 2021 hides earlier histories, which pulls the entry contrasts towards zero. Some abstracts and titles
are final versions written after the decision; code links were removed for discovery papers, but titles could not be
corrected. Withdrawn papers are left out of the acceptance tests. Finally, the hindsight comparison is a bracket,
because the topic maps differ in the papers they are fitted on as well as in their time window.

\section{Future work}
\label{sec:future}

The most direct next step is to separate status from continuity. If the record-side associations reflect status, they
should be stronger where a preprint identifies the authors before review, which is how \citet{ibrahim2026bias} found
prestige to operate. Running this test would require a new declaration before any outcome is read, as the protocol of
this report requires.

The labels also need a human check that does not depend on one person. Two coders labelling the same random sample
independently would give agreement between humans and between humans and the LLM, and would show whether the
borderline cases change the entry results. A related question is what newcomers' method papers contain: whether they
propose new ideas or adapt existing ones to a new topic.

Several open questions could be answered with more data. Whether entrants flock to hot topics remains unsettled, and a
registered held-out test with enough power could settle it. Linking submissions to their later citations would show
whether the work that is accepted more often is also the work that has lasting impact, as the persistence results of
\citet{bao2025persistence} suggest it may not be. Withdrawn papers and the text of reviews could show whether the gate
also selects through withdrawal and what reviewers object to in papers that move far from their authors' earlier work.
Data on compute or institutions could show whether the concentration of measurement papers among incumbents reflects
resources other than team size.

Finally, the design could be repeated at other ML and NLP venues, such as NeurIPS or ACL, to see whether the findings
are specific to ICLR. Refitting topic maps on early years and using predictors measured only from the past could also
become a routine check for claims about research trends.

\section{Closing remarks}

Machine learning is easy to describe as a field that chases fashions, and its keyword waves and the rapid spread of
foundation models seem to confirm it. When the herding account is tested where it makes distinct predictions, however,
it finds little support at ICLR. Newcomers to a topic bring methods rather than follow-on work, and acceptance shows
no detectable association with following the crowd. What acceptance is associated with, instead, is staying close to
the team's own earlier work. Whether that reflects expertise, status or the value of continuing a line of
work is the main question this report leaves open.
